\chapter{Methodology}
\label{chap:method}

\section{Overview: A Progression of Approaches}
We investigated three method families of increasing sophistication. Approach~A
(direct prompting) established baselines and exposed failure modes. Approach~B
(multi-persona voting) addressed output bias and single-pass variance through prompt
diversity and majority voting on a local model. Approach~C (VERA-lite) introduced
guiding-question optimisation, a blind evaluation protocol, alternative input
representations, and ensembling, evaluated with repeated runs.

\section{Datasets}

\subsection{IPAD and the Motivation for a Custom Dataset}
Our early experiments (Approaches~A and~B) used the IPAD industrial benchmark, in
particular its \textbf{R01 push-button} scene, in which a small cylindrical push-button
travels on a belt and the dominant anomaly is a reversed orientation (the red cap on the
wrong side). IPAD has two drawbacks for our purposes. First, its frames are low
resolution ($256\times256$), which limits the visibility of the fine-grained cues we care
about, such as a missing cap or a subtle colour difference on a small part. Second, and
more decisively, the benchmark does not let us control the specific anomaly types we wanted
to study, nor does it provide a small, cleanly labelled set tailored to a focused,
per-anomaly-type quantitative analysis. For these reasons we recorded our own dataset
under controlled conditions, with a fixed camera and a documented protocol, which also let
us define the anomaly taxonomy and obtain reliable ground-truth labels.

\subsection{The Pen-Conveyor Dataset}
The \textbf{pen-conveyor} scene is the purpose-built dataset used for the main (VERA-lite)
experiments: a white board slides left-to-right on a linear rail carrying three identical
pens. It contains $13$ clips ($3$ normal and $10$ anomalous) spanning seven anomaly types:
missing cap, wrong orientation, wrong colour, missing pen, belt stop/freeze, unexpected
(intra-object) pen movement, and a combined case. Clips were recorded at 4K and downscaled
to $720\times1280$ for API use; ground-truth labels follow the documented recording
protocol. Figure~\ref{fig:dataset} shows a representative frame for the normal class and
for each anomaly type.

\begin{figure}[!htbp]
    \centering
    \includegraphics[width=\textwidth]{fig_dataset.jpg}
    \caption{\label{fig:dataset}The pen-conveyor dataset: a representative frame for the
    normal class and for each anomaly type (missing cap, wrong orientation, wrong colour,
    missing pen, belt stop, and unexpected intra-object movement).}
\end{figure}

\section{Approach A: Direct Prompting}
We first prompted an API-served Qwen-VL model directly, in four modes: (i) zero-shot
single frame, (ii) few-shot with a normal reference frame, (iii) a multi-frame sequence
sent as separate images with a sliding window, and (iv) a frame-to-video clip. Three
prompt styles were compared (plain, chain-of-thought~\cite{cot}, and attribute-based). This phase
showed that chain-of-thought / step-wise prompts are the most consistent and that the
model tends to default to ``normal'' under terse output formats.

\section{Approach B: Multi-Persona Majority Voting}
\label{sec:persona}
To reduce single-pass variance and the ``normal''-bias, we run the same clip through
several \emph{inspector personas} and decide by majority vote. This runs on a locally
deployed open-weight model (\texttt{Qwen3-VL-8B-Instruct}), so data never leaves the
premises.

\subsection{Voting Procedure}
Twenty personas were authored (line operators, mechanical/quality/vision/robotics
engineers, auditors, lean and Toyota experts, conservative and aggressive auditors,
etc.). All share the \emph{same} anomaly rules but differ in voice, threshold and output
format. A voting subset of five personas plus one distilled \emph{consensus} prompt
forms six voters; verdicts are parsed to NORMAL/ANOMALY/UNCLEAR and combined by majority
vote, with confidence given by the winning fraction.

\subsection{The Three Persona Types (Ablation)}
The personas were studied at three checklist breadths
(Table~\ref{tab:persona-types}): every voter checks a single focus area
(\emph{specialized}), every voter checks the full A/B/C checklist (\emph{complete}), or
every voter answers only the orientation question (\emph{orientation}).

\begin{table}[!htbp]
\centering
\caption{\label{tab:persona-types}The three persona ablations.}
\begin{tabular}{l|l|l}
\hline
Type & What each voter checks & Rationale \\
\hline
specialized & A single focus area & Committee of specialists \\
complete    & Full A/B/C checklist & Redundancy (full inspection each) \\
orientation & Orientation only     & Minimal / primary-anomaly focus \\
\hline
\end{tabular}
\end{table}

\subsection{The Articulate-First Finding}
Compact YES/NO output formats produced false negatives: the model assigned ``normal''
without actually examining the frames. Forcing an \emph{articulate-first} format --- the
model first describes the entry/middle/exit frames, then derives each check, then gives
the verdict --- suppressed this bias. This mirrors the question-decomposition principle
of Approach~C: making the model verbalise its observations before deciding improves
reliability.

\section{Approach C: VERA-lite}
\label{sec:vera}
VERA-lite keeps the model frozen and improves the \emph{guiding questions}.

\subsection{Pipeline}
A \emph{learner} VLM answers the guiding questions for a clip and outputs a JSON verdict
(NORMAL/ANOMALY), an anomaly type, a score and a confidence. An \emph{optimiser} VLM
inspects the learner's mistakes and proposes improved questions. One iteration runs the
learner over all clips, then the optimiser.

The final (v3) guiding questions used for the pen-conveyor scene are:
\begin{enumerate}
  \item When the board is fully visible, does it hold exactly three pens arranged in a
        single row?
  \item Are all three pens aligned identically, all pointing the same way, including any
        pen near the left or right frame edge as it enters or exits?
  \item Do all three pens share identical barrel and cap colours, with no mismatch?
  \item Is every pen fully capped with no exposed writing tip, and are all three pens
        equal in length?
  \item Does the board advance steadily from left to right at every moment, never
        halting, freezing, or pausing even briefly (1--3\,s) at any point?
  \item Do all pens stay fixed in their positions on the board, with no pen shifting,
        sliding, or moving on its own relative to the board mid-clip?
\end{enumerate}
Each question targets one anomaly family (count, orientation, colour, cap/completeness,
belt motion, and intra-object movement, respectively).

\subsection{Blind Evaluation Protocol}
The original pipeline leaked the expected label into the learner prompt. We made
\emph{blind} mode the default: the learner never sees the ground-truth label, which is
used only for scoring. All reported accuracies are blind. We added accuracy, precision,
recall, $F_1$ and a confusion matrix to the run summary, plus retry/back-off for
transient API errors.

\subsection{Input Representations}
Because a single video call samples frames sparsely, we compared representations:
the native \textbf{video} call (with a controllable sampling rate \texttt{fps}); a
\textbf{dense-frame} mode that sends $N$ time-ordered frames as separate numbered images
with a relative-spacing prompt, restricted to a central window to avoid entry/exit
artifacts; and a \textbf{temporal grid} that tiles frames into a single montage image.

\subsection{Reasoning-Budget Control and Ensembling}
The reasoning model can exhaust its output budget on internal deliberation and return an
empty answer; a \texttt{thinking\_budget} caps the reasoning and forces a final answer,
which also controls latency. Finally, a two-view \textbf{OR-ensemble} combines the video
and dense-frame verdicts (ANOMALY if either flags), exploiting their disjoint blind
spots.

\section{Evaluation Metrics}
We report clip-level binary accuracy and, treating ANOMALY as the positive class,
precision, recall and $F_1$, together with the confusion counts (TP/TN/FP/FN). Because
VLM outputs are non-deterministic, key configurations are run multiple times and we
report the mean and standard deviation. Latency is reported as seconds per clip.
